% -*- coding: utf-8 -*-
% !TEX program = xelatex
\documentclass[plain,noanswer,medmath]{../../template/jnuexam}
\usepackage{../../template/kysx}

\begin{document}


\examtitle{1988}{4}


\makepart{填空题}{本题满分12分，每空1分}


\begin{problem}
已知函数$f(x)=\int_0^x\e^{-\frac{1}{2}t^2}\dt$，$-\infty<x<+\infty$，则\par
(a) $f'(x)=$ \fillin{}；\par
(b) $f(x)$的单调性：\fillin{}；\par
(c) $f(x)$的奇偶性：\fillin{}；\par
(d) $f(x)$的图形的拐点：\fillin{}；\par
(e) $f(x)$图形的凹凸性：\fillin{}，\fillin{}；\par
(f) $f(x)$图形的水平渐近线：\fillin{}，\fillin{}．
\end{problem}


\begin{problem}
$\begin{vmatrix}
  1 & 1 & 1 & 0 \\
  1 & 1 & 0 & 1 \\
  1 & 0 & 1 & 1 \\
  0 & 1 & 1 & 1
\end{vmatrix}=$ \fillin{}．
\end{problem}


\begin{problem}
$\begin{pmatrix}
  0 & 0 & 0 & 1 \\
  0 & 0 & 1 & 0 \\
  0 & 1 & 0 & 0 \\
  1 & 0 & 0 & 0
\end{pmatrix}^{-1}=$ \fillin{}．
\end{problem}


\begin{problem}
假设$P(A)=0.4$，$P(A\cup B)=0.7$，那么\par
(a) 若$A$与$B$互不相容，则$P(B)=$ \fillin{}；\par
(b) 若$A$与$B$相互独立，则$P(B)=$ \fillin{}．
\end{problem}


\makepart{判断题}{本题满分10分，每小题2分}
%本题满分10分；每小题回答正确得2分，回答错误得-1分，不回答得0分；全题最低得0分

\begin{problem}
若极限$\lim_{x \to x_0} f(x)$与$\lim_{x \to x_0} f(x)g(x)$都存在，
则极限$\lim_{x \to x_0} g(x)$必存在．\tickout{}
\end{problem}


\begin{problem}
若$x_0$是函数$f(x)$的极值点，则必有$f'(x_0) = 0$． \tickout{}
\end{problem}


\begin{problem}
等式$\int_0^a f(x)\dx = -\int_0^a f(a-x)\dx$对任何实数$a$都成立．\tickout{}
\end{problem}


\begin{problem}
若$A$和$B$都是$n$阶非零方阵，且$AB=0$，则$A$的秩必小于$n$． \tickout{}
\end{problem}


\begin{problem}
若事件$A$, $B$, $C$满足等式$A\cup C = B \cup C$，则$A=B$． \tickout{}
\end{problem}


\makepart{计算题}{本题满分16分，每小题4分}


\begin{problem}
求极限$\lim_{x \to 1} \frac{x^x - 1}{x\ln x}$．
\end{problem}


\begin{problem}
已知$u + \e^u = xy$，求$\frac{\pd^2 u}{\pdx\pd y}$．
\end{problem}


\begin{problem}
求定积分$\int_0^3\frac{\dx}{\sqrt{x}(1 + x)}$．
\end{problem}


\begin{problem}
求二重积分$\int_0^{\frac{\pi}{6}} \dy \int_y^{\frac{\pi}{6}} \frac{\cos x}{x}\dx$．
\end{problem}


\makepart{解答题}{本题满分6分，每小题3分}


\begin{problem}
讨论级数$\sum_{n = 1}^{\infty}\frac{(n + 1)!}{n^{(n + 1)}}$的敛散性．
\end{problem}


\begin{problem}
已知级数$\sum_{n = 1}^{\infty}a_n^2$和$\sum_{n = 1}^{\infty}b_n^2$都收敛，
试证明级数$\sum_{n = 1}^{\infty}a_n b_n$绝对收敛，
\end{problem}


\makepart{}{本题满分8分}

\begin{problem}
已知某商品的需求量$D$和供给量$S$都是价格$p$的函数： 
$$D = D(p) = \frac{a}{p^2},\qquad S = S(p) = bp,$$
其中$a > 0$和$b > 0$为常数；价格$p$是时间$t$的函数且满足方程
$$\frac{\d p}{\dt}=k[D(p)-S(p)]\qquad \text{（$k$为正的常数）}.$$
假设当$t = 0$时价格为$1$，试求
\begin{enumerate}
\item 需求量等于供给量时的均衡价格$p_e$；
\item 价格函数$p(t)$；
\item 极限$\lim_{t \to +\infty} p(t)$．
\end{enumerate}
\end{problem}


\makepart{}{本题满分8分}

\begin{problem}
在曲线$y = x^2$（$x > 0$）上某点$A$处作一切线，
使之与曲线以及$x$轴所围图形的面积为$\frac{1}{12}$．试求：
\begin{enumerate}
\item 切点$A$的坐标；
\item 过切点$A$的切线方程；
\item 由上述所围平面图形绕$x$轴旋转一周所成旋转体的体积．
\end{enumerate}
\end{problem}


\makepart{}{本题满分8分}

\begin{problem}
已给线性方程组$\begin{cases}
  x_1 + x_2 + 2x_3 + 3x_4 = 1, \\
  x_1 + 3x_2 + 6x_3 + x_4 = 3, \\
  3x_1 - x_2 - k_1x_3 + 15x_4 = 3, \\
  x_1 - 5x_2 - 10x_3 + 12x_4 = k_2.
\end{cases}$问$k_1$和$k_2$各取何值时，方程组无解？有唯一解？有无穷多解？
在方程组有无穷多组解的情形下，试求出一\text{般解}．
\end{problem}


\makepart{}{本题满分7分}

\begin{problem}
已知向量组$\alpha_1,\alpha_2, \cdots ,\alpha_s$ ($s \ge 2$)线性无关．设
$$\beta_1 = \alpha_1 + \alpha_2,\quad \beta_2 = \alpha_2 + \alpha_3,\quad \cdots,\quad
  \beta_{s - 1} = \alpha_{s - 1} + \alpha_s,\quad \beta_s = \alpha_s + \alpha_1.$$
试讨论向量组$\beta_1,\beta_2, \cdots ,\beta_s$的线性相关性．
\end{problem}


\makepart{}{本题满分6分}

\begin{problem}
设$A$是三阶方阵，$A^*$是$A$的伴随矩阵，$A$的行列式$|A| = \frac{1}{2}$．
求行列式$\left|(3A)^{-1} - 2A^*\right|$的值．
\end{problem}


\makepart{}{本题满分7分}

\begin{problem}
玻璃杯成箱出售，每箱$20$只，假设各箱含$0$, $1$, $2$只残次品的概率相应为$0.8$, $0.1$ 和$0.1$．
一顾客欲购一箱玻璃杯，在购买时，售货员随意取一箱，顾客开箱随机地察看$4$只：
若无残次品，则买下该箱玻璃杯，否则退回．试求：
\begin{enumerate}
\item 顾客买下该箱的概率$\alpha$；
\item 在顾客买下的一箱中，确实没有残次品的概率$\beta$．
\end{enumerate}
\end{problem}


\makepart{}{本题满分6分}

\begin{problem}
某保险公司多年的统计资料表明，在索赔户中被盗索赔户占$20\%$，
以$X$表示在随意抽查的$100$个索赔户中因被盗向保险公司索赔的户数．
\begin{enumerate}
\item 写出$X$的概率分布；
\item 利用棣莫佛—拉普拉斯定理，求被盗索赔户不少于$14$户且不多于$30$户的概率的近似值．
\end{enumerate}
\centerline{[附表]设$\Phi(x)$是标准正态分布函数．}%
\noindent\centerline{$\begin{array}{|c|ccccccc|}
\hline
        x &     0 &   0.5 &   1.0 &   1.5 &   2.0 &   2.5 & 3.0 \\
\hline
  \Phi(x) & 0.500 & 0.692 & 0.841 & 0.933 & 0.977 & 0.994 & 0.999 \\
\hline
\end{array}$}
\end{problem}


\makepart{}{本题满分6分}

\begin{problem}
假设随机变量$X$在区间$(1,2)$上服从均匀分布，试求随机变量$Y = \e^{2X}$的概率密度$f_Y(y)$．
\end{problem}


\end{document}
